% texte.tex — f162v, vv. 1–42 (conversion initiale des travaux antérieurs)
\rubrique{ΠΡΟΛΟΓΟΣ}{Prologue}
\vers[n=1]{ΜΝ.}{Mnésiloque}{ὦ Ζεῦ χελιδὼν ἆρά ποτε φανήσεται;}{Ô Zeus, est-ce qu'à la fin une hirondelle paraîtra ?}
\vers[n=2]{}{}{ἀπολεῖ μ᾽ ἀλοῶν ἅνθρωπος ἐξ ἑωθινοῦ.}{Cet homme me tuera de me faire tribuler dès le matin.}
\vers[n=3]{}{}{οἷόν τε, πρὶν τὸν σπλῆνα κομιδῇ μ᾽ ἐκβαλεῖν.}{Y a-t-il moyen, avant que perdre tout de bon la rate,}
\vers[n=4]{}{}{παρὰ σοῦ πυθέσθαι ποῖ μ᾽ ἄγεις ωὖριπίδη;}{de savoir de toi où tu me mènes, Euripide ?}
\vers[n=5]{ΕΥ.}{Euripide}{ἀλλ᾽ οὐκ ἀκούειν δεῖ σε πάνθ᾽ ὅσ᾽ αὐτίκα}{Eh ! il n'est pas besoin que tu entendes tout ce que}
\vers[n=6]{}{}{ὄψει παρεστώς.}{tout à l'heure tu verras de tes yeux.}
\vers[suite]{ΜΝ.}{Mnésiloque}{πῶς λέγεις; αὖθις φράσον.}{Comment dis-tu ? Répète.}
\vers[n=7]{}{}{οὐ δεῖ μ᾽ ἀκούειν;}{Il n'est pas besoin que j'entende.}
\vers[suite]{ΕΥ.}{Euripide}{οὐχ ἅ γ᾽ ἂν μέλλῃς ὁρᾶν.}{Non, ce que tu vas voir.}
\vers[n=8]{ΜΝ.}{Mnésiloque}{οὐδ᾽ ἆρ᾽ ὁρᾶν δεῖ μ᾽;}{Ni donc que je voie...}
\vers[suite]{ΕΥ.}{Euripide}{οὐχ ἅ γ᾽ ἂν ἀκούειν δέῃ.}{Non, ce qu'il faut que tu entendes.}
\vers[n=9]{ΜΝ.}{Mnésiloque}{πῶς μοι παραινεῖς; δεξιῶς μέντοι λέγεις.}{Singulier conseil ! Ma foi, c'est finement dit.}
\vers[n=10]{}{}{οὐ φῂς σὺ χρῆναί μ᾽ οὔτ᾽ ἀκούειν οὔθ᾽ ὁρᾶν.}{Tu prétends que je ne dois ni entendre ni voir.}
\vers[n=11]{ΕΥ.}{Euripide}{χωρὶς γὰρ αὐτοῖν ἑκατέρου 'στὶν ἡ φύσις.}{C'est que distincte est l'essence de chacune de ces choses.}
\vers[n=12]{ΜΝ.}{Mnésiloque}{τοῦ μήτ᾽ ἀκούειν μήθ᾽ ὁρᾶν;}{Ne pas entendre et ne pas voir ?}
\vers[suite]{ΕΥ.}{Euripide}{εὖ ἴσθ᾽ ὅτι.}{Tiens-le pour assuré.}
\vers[n=13]{ΜΝ.}{Mnésiloque}{πῶς χωρίς;}{Comment distincte ?}
\vers[suite]{ΕΥ.}{Euripide}{οὕτω ταῦτα διεκρίθη τότε.}{Voici comment elles furent départies à l'origine. L'Éther, quand d'abord il se dégagea}
\vers[n=14]{}{}{αἰθὴρ γὰρ ὅτε τὰ πρῶτα διεχωρίζετο}{et procréa}
\vers[n=15]{}{}{καὶ ζᾦ᾽ ἐν αὑτῷ ξυνετέκνου κινούμενα,}{en soi-même des animaux doués de mouvement,}
\vers[n=16]{}{}{ᾧ μὲν βλέπειν χρὴ πρῶτ᾽ ἐμηχανήσατο}{commença par fabriquer pour qui doit voir}
\vers[n=17]{}{}{ὀφθαλμὸν ἀντίμιμον ἡλίου τροχῷ,}{un œil analogique de l'orbe solaire}
\vers[n=18]{}{}{ἀκοὴν δὲ χοάνης ὦτα διετετρήνατο.}{et fora une oreille à l'image des godets d'un moule.}
\vers[n=19]{ΜΝ.}{Mnésiloque}{διὰ τὴν χοάνην οὖν μήτ᾽ ἀκούω μήθ᾽ ὁρῶ;}{Et à cause de ce moule je ne dois ni entendre ni voir ?}
\vers[n=20]{}{}{νὴ τὸν Δί᾽ ἥδομαί γε τουτὶ προσμαθών.}{Par Zeus, je suis ravi d'apprendre encore cela.}
\vers[n=21]{}{}{οἷόν γέ που 'στιν αἱ σοφαὶ ξυνουσίαι.}{Quelle chose tout de même, que les doctes accointances !}
\vers[n=22]{ΕΥ.}{Euripide}{πόλλ᾽ ἂν μάθοις τοιαῦτα παρ᾽ ἐμοῦ.}{Auprès de moi, tu en apprendras bien d'autres du même genre.}
\vers[suite]{ΜΝ.}{Mnésiloque}{πῶς ἂν οὖν}{Ah ! que ne puis-je}
\vers[n=24]{}{}{ἔτι προσμάθοιμι χωλὸς εἶναι τὼ σκέλει;}{en outre apprendre aussi à être perclus des jambes ?}
\vers[n=25]{ΕΥ.}{Euripide}{βάδιζε δευρὶ καὶ πρόσεχε τὸν νοῦν.}{Avance ici et prête-moi ton attention.}
\vers[suite]{ΜΝ.}{Mnésiloque}{ἰδού.}{J'y suis.}
\vers[n=26]{ΕΥ.}{Euripide}{ὁρᾷς τὸ θύριον τοῦτο;}{Tu vois cette petite porte ?}
\vers[suite]{ΜΝ.}{Mnésiloque}{νὴ τὸν Ἡρακλέα}{Oui, par Héraclès,}
\vers[n=27]{}{}{οἶμαί γε.}{à ce que je crois, du moins.}
\vers[suite]{ΕΥ.}{Euripide}{σίγα νυν.}{Maintenant, motus.}
\vers[suite]{ΜΝ.}{Mnésiloque}{σιωπῶ τὸ θύριον.}{Je ne dis mot de la petite porte.}
\vers[n=28]{ΕΥ.}{Euripide}{ἄκου᾽.}{Écoute !}
\vers[suite]{ΜΝ.}{Mnésiloque}{ἀκούω καὶ σιωπῶ τὸ θύριον.}{J'écoute et ne dis mot de la petite porte.}
\vers[n=29]{ΕΥ.}{Euripide}{ἐνταῦθ᾽ Ἀγάθων ὁ κλεινὸς οἰκῶν τυγχάνει}{C'est là que se trouve demeurer l'illustre Agathon,}
\vers[n=30]{}{}{ὁ τραγῳδοποιός.}{le poète tragique.}
\vers[suite]{ΜΝ.}{Mnésiloque}{ποῖος οὗτος Ἁγάθων;}{Qui est-ce, cet Agathon ?}
\vers[n=31]{}{}{ἔστιν τις ἀγαθῶν; μῶν ὁ μέλας ὁ καρτερός;}{Est-ce l'un de ces honnêtes hommes... celui au teint hâlé, plein d'allant ?}
\vers[n=32]{ΕΥ.}{Euripide}{οὔκ, ἀλλ᾽ ἕτερός τις· οὐχ ἑόρακας πώποτε;}{Non, c'est un autre. Tu ne l'as jamais vu ?}
\vers[n=33]{ΜΝ.}{Mnésiloque}{μῶν ὁ δασυπώγων;}{L'homme à la barbe touffue ?}
\vers[suite]{ΕΥ.}{Euripide}{οὐχ ἑόρακας πώποτε.}{Tu ne l'as jamais vu !..}
\vers[n=34]{ΜΝ.}{Mnésiloque}{μὰ τὸν Δί᾽ οὔτοι γ᾽ ὥστε κἀμέ γ᾽ εἰδέναι.}{Non, par Zeus, pas du moins que je sache.}
\vers[n=35]{ΕΥ.}{Euripide}{καὶ μὴν †βεβίνηκας† σύ γ᾽, ἀλλ᾽ οὐκ οἶσθ᾽ ἴσως.}{À coup sûr tu l'as enfourché, mais peut-être ne t'en souvient-il pas.}
\vers[n=36]{}{}{ἀλλ᾽ ἐκποδὼν πτήξωμεν, ὡς ἐξέρχεται}{Mais allons nous tapir à l'écart. Car sort de chez lui}
\vers[n=37]{}{}{θεράπων τις αὐτοῦ πῦρ ἔχων καὶ μυρρίνας}{un serviteur avec du feu et des myrtes.}
\vers[n=38]{}{}{προθυσόμενος ἔοικε τῆς ποιήσεως.}{Il va sacrifier, ce semble, en faveur du succès de sa poésie.}
\vers[n=39, indent=3em, findent=1.7999999999999998em, metre={an.}]{ΘΕ.}{Serviteur}{εὔφημος πᾶς ἔστω λαός,}{Que le peuple entier se recueille}
\vers[n=40, indent=3em, findent=1.7999999999999998em]{}{}{στόμα συγκλῄσας· ἐπιδημεῖ γὰρ}{et ferme la bouche, car est venu séjourner}
\vers[n=41, indent=3em, findent=1.7999999999999998em]{}{}{θίασος Μουσῶν ἔνδον μελάθρων}{un thiase de muses dans la demeure}
\vers[n=42, indent=3em, findent=1.7999999999999998em]{}{}{τῶν δεσποσύνων μελοποιῶν.}{de mon maître pour composer des chants.}
